\documentclass[10pt,twoside]{article}

% Encoding and font
\usepackage[T1]{fontenc}           % Modern font encoding
\usepackage{lmodern}               % Scalable vector fonts (replaces Computer Modern)
\renewcommand*\familydefault{\sfdefault}  % Sans-serif as default

% Layout and math
\usepackage[left=2.5cm, top=2.5cm, bottom=2.5cm, right=3cm]{geometry}
\usepackage{graphicx, amssymb, amsmath}
\usepackage{fancyhdr}
\usepackage{setspace}
\usepackage{lineno}
\usepackage{longtable}
\usepackage{pdflscape}     % For landscape pages
\usepackage[table]{xcolor}
\usepackage{url}

% Figures
\usepackage{subcaption}            % Subfigure support

% Author and affiliation support
\usepackage{authblk}

% Header/footer customization
\pagestyle{fancy}
\fancyhead{}
\fancyfoot{}
\fancyhead[RO,LE]{Supporting Information}
\fancyfoot[RO,LE]{Page-\thepage}

% Title and authorship
\title{
  Supporting Information - 3 - Software Setup: \\
  \textit{San Jos\'{e}: A Reproducible, Bias-Aware Reference Layer for timsTOF Data on PRIDE}
}

\author[1,2]{David Teschner}
% TODO: add co-authors
\author[1,2]{Andreas Hildebrandt}
\author[3,4,5]{Stefan Tenzer}

\affil[1]{Institute of Computer Science, Johannes Gutenberg University, 55128 Mainz, Germany}
\affil[2]{Institute for Quantitative and Computational Biosciences (IQCB), Johannes Gutenberg University, 55128 Mainz, Germany}
\affil[3]{Helmholtz Institute for Translational Oncology, Mainz, Germany}
\affil[4]{German Cancer Research Center (DKFZ), 69120 Heidelberg, Germany}
\affil[5]{University Medical Center, Johannes Gutenberg University, 55131 Mainz, Germany}

\date{\today}

\begin{document}

\onehalfspacing
\maketitle

\section{Installation}

\subsection{Python Environment}

San Jos\'{e} requires Python~3.11 or later. We recommend creating a dedicated virtual environment:

\begin{verbatim}
python3.12 -m venv .venv
source .venv/bin/activate
pip install -r requirements.txt
\end{verbatim}

The rustims framework (imspy-core, imspy-connector) requires a Rust toolchain for compilation. If pre-built wheels are not available, install Rust via \url{https://rustup.rs/} before running pip install.

\begin{verbatim}
# Build imspy from source (if needed)
./scripts/setup_imspy.sh
\end{verbatim}

\subsection{Search Engines}

Three external search engines must be installed separately:

\begin{longtable}{|l|l|l|p{6cm}|}
\hline
\rowcolor{gray!20}
\textbf{Engine} & \textbf{Version} & \textbf{License} & \textbf{Installation} \\
\hline
FragPipe & 23+ & Academic & Download from \url{https://fragpipe.nesvilab.org/}. Requires academic license agreement. \\
\hline
DIA-NN & 1.9+ & Academic & Download from \url{https://github.com/vdemichev/DiaNN} \\
\hline
Sage & 0.14+ & MIT & Download from \url{https://github.com/lazear/sage/releases} or build from source with Rust \\
\hline
\end{longtable}

\section{Configuration}

All paths and parameters are configured in \texttt{config/config.yaml}:

\begin{verbatim}
# Dataset selection
datasets:
  - PXD019086
  - PXD046675

# Organism FASTA databases
organisms:
  human:
    taxon_id: 9606
    proteome_id: "UP000005640"
    local_fasta: "/path/to/human.fasta"
  yeast:
    taxon_id: 559292
    proteome_id: "UP000002311"

# Search engine paths (user must provide)
fragpipe:
  path: "/path/to/fragpipe"
  workflow: "LFQ-MBR"
  threads: 64
  memory_gb: 128

diann:
  path: "/path/to/diann"
  threads: 64
  qvalue: 1.0    # Report ALL results

sage:
  path: "/path/to/sage"
  threads: 64

# Extraction parameters
extraction:
  mz_tolerance_ppm: 20
  rt_tolerance_sec: 15
  batch_size: 10000

# Test mode (for quick validation)
test_mode:
  enabled: false
  max_files: 3
\end{verbatim}

\section{Usage}

\subsection{Running the Pipeline}

\begin{verbatim}
# Full pipeline for a PRIDE accession
.venv/bin/python runner/run_dataset.py PXD019086 \
    --config config/config.yaml

# Test mode (limited files for quick validation)
.venv/bin/python runner/run_dataset.py PXD019086 \
    --test-mode --max-files 1

# Use local data (skip download)
.venv/bin/python runner/run_dataset.py PXD019086 \
    --local-data data/raw/PXD019086

# Resume from checkpoint after interruption
.venv/bin/python runner/run_dataset.py PXD019086 --resume

# Run specific steps only
.venv/bin/python runner/run_dataset.py PXD019086 --steps 4 5

# Create distributable archive
.venv/bin/python runner/run_dataset.py PXD019086 \
    --package --package-version 1.0
\end{verbatim}

\subsection{Running the Dashboard}

\begin{verbatim}
# Backend (single dataset)
.venv/bin/python dashboard/backend/main.py \
    --store data/merged/PXD019086/precursor_store.parquet \
    --port 8000

# Backend (collection mode, multiple datasets)
.venv/bin/python dashboard/backend/main.py \
    --collection /path/to/san_jose_collection/ \
    --port 8000

# Frontend (serve pre-built files)
cd dashboard/frontend/dist && python3 -m http.server 5173
\end{verbatim}

Open \url{http://localhost:5173} in a browser to access the precursor browser.

\subsection{Adding Datasets to a Collection}

\begin{verbatim}
python scripts/add_to_collection.py \
    --archive data/packages/PXD019086_v1.0.zip \
    --collection /path/to/collection \
    --study meier_2021_ccs
\end{verbatim}

\subsection{HPC Execution (SLURM)}

For HPC environments with Singularity:

\begin{verbatim}
# Using Snakemake with SLURM profile
snakemake add_all_datasets --profile profiles/mogon2

# Or direct SLURM submission
sbatch scripts/slurm_run_dataset.sh PXD019086
\end{verbatim}

\section{Output Structure}

After processing, the following directory structure is produced:

\begin{verbatim}
data/
  raw/{accession}/                  # Step 1: Bruker .d folders
  processed/{accession}/            # Steps 2-3
    fragpipe_output/                #   FragPipe results
    diann/                          #   DIA-NN results
    sage/                           #   Sage results
    consensus/stratified/           #   Precursors by agreement tier
    precursor_index.parquet         #   Unified index
  extracted/{accession}/            # Step 4: Raw 4D signal
    {raw_file}.d/blobs.bin          #   Raw 4D data blobs
  merged/{accession}/               # Step 5: Final output
    precursor_store.parquet         #   IDs + raw features joined
    manifest.json                   #   QC metrics
  packages/                         # Step 6: Archives
    {accession}_v{version}.zip
  checkpoints/{accession}/          # Pipeline state
    state.json
\end{verbatim}

\end{document}
